% Generated draft: requires mathematical and source review.
\documentclass[11pt]{article}
\usepackage[margin=25mm]{geometry}
\usepackage{fontspec}
\setmainfont{Latin Modern Roman}
\usepackage{amsmath,amssymb,mathtools,unicode-math}
\setmathfont{Latin Modern Math}
\usepackage{enumitem,needspace,xurl,hyperref}
\setlength{\parindent}{0pt}
\setlength{\parskip}{5pt}
\setlength{\emergencystretch}{3em}
\newcommand{\N}{\mathbb{N}}
\newcommand{\Z}{\mathbb{Z}}
\newcommand{\Q}{\mathbb{Q}}
\newcommand{\R}{\mathbb{R}}
\newcommand{\C}{\mathbb{C}}
\newcommand{\F}{\mathbb{F}}
\newcommand{\A}{\mathbb{A}}
\newcommand{\PP}{\mathbb{P}}
\newcommand{\E}{\mathbb{E}}
\newcommand{\eps}{\varepsilon}
\newcommand{\dd}{\,\mathrm{d}}
\begin{document}
\section*{UnsolvedMath ID: 30004457}
\textbf{Source problem code:} catalogue entry
\subsection{Definitions and mathematical statement}
Let \(G\) and \(H\) be finitely generated groups, equipped with word metrics \(d_G\) and \(d_H\) arising from finite symmetric generating sets. A quasi-isometry from \(G\) to \(H\) is a map \(f:G\to H\) for which there exist constants \(L\geq 1\) and \(C\geq 0\) such that
\[
L^{-1}d_G(g,g')-C\leq d_H(f(g),f(g'))\leq Ld_G(g,g')+C
\qquad(g,g'\in G),
\]
and every point of \(H\) lies within distance \(C\) of \(f(G)\). The groups are quasi-isometric if such a map exists; this property is independent of the chosen finite generating sets.

A geodesic metric space \(X\) is Gromov hyperbolic if there exists \(\delta\geq 0\) such that every side of each geodesic triangle lies in the closed \(\delta\)-neighborhood of the other two sides. An isometric action \(G\curvearrowright X\) is acylindrical if, for every \(\varepsilon>0\), there exist \(R\geq 0\) and \(N\in\N\) such that, for all \(x,y\in X\) with \(d_X(x,y)\geq R\),
\[
\#\{g\in G: d_X(x,gx)\leq\varepsilon\text{ and }d_X(y,gy)\leq\varepsilon\}\leq N.
\]
The action is non-elementary if its limit set in the Gromov boundary \(\partial X\) has at least three points. Here the limit set consists of boundary limits of sequences in an orbit \(Gx\). A group is acylindrically hyperbolic if it admits a non-elementary acylindrical isometric action on a Gromov hyperbolic geodesic metric space. No properness or cocompactness assumption on this action is imposed.

\textbf{Question.} For every pair of finitely generated groups \(G,H\), does
\[
G\text{ is acylindrically hyperbolic and }G\text{ is quasi-isometric to }H
\quad\Longrightarrow\quad
H\text{ is acylindrically hyperbolic}
\]
hold?

The source also asks, more broadly, whether the class of finitely generated acylindrically hyperbolic groups is quasi-isometrically rigid. If this means that the class is closed under quasi-isometry among finitely generated groups, it is exactly the displayed question. The record specifies no stronger rigidity conclusion or equivalence relation, so a distinct stronger mathematical assertion cannot be extracted from that wording.

The supplied background mentions a special result preserving acylindrical hyperbolicity and substantial partial criteria, but gives neither their hypotheses nor their conclusions. These cannot be incorporated as precise auxiliary statements. The record labels the problem open; that label and its literature assessment are supplied curation, not independently verified mathematical conclusions.
\subsection{Short English statement}
Must every finitely generated group quasi-isometric to a finitely generated acylindrically hyperbolic group also be acylindrically hyperbolic? The additional phrase ``quasi-isometrically rigid'' has no specified stronger meaning in the supplied record. The supplied open status requires independent review.
\subsection{Sources}
\begin{itemize}
\item Local registry: \texttt{data/unsolvedmath/problems.json}; supplementary registry: \texttt{catalog/problems.json}. Supplied permanent UnsolvedMath ID: \(30004457\). Exact source problem code: \texttt{catalogue entry}. Collection: Oberwolfach Reports Open Problems. The supplied proposed year is 2020; no individual proposer is named.
\item Supplied source citation: \emph{Geometric Structures in Group Theory} (2021), Oberwolfach Reports, pp.\ 877--918, DOI 10.4171/owr/2020/16. \url{https://doi.org/10.4171/owr/2020/16}. The background also labels this source ``Oberwolfach Report 16/2020'' and describes the reports as documenting mathematical workshops and their open-problem discussions.
\item Supplied literature source: Hughes--Martinez-Pedroza (2023). \url{https://doi.org/10.4153/S0008439523000012}. No article title or precise theorem reference is supplied.
\item Supplied literature source: Oberwolfach Report 16/2020. \url{https://doi.org/10.4171/owr/2020/16}.
\item Supplied literature assessment, checked on \texttt{2026-08-22}: the general quasi-isometry-invariance question remains open, while a special preservation result and substantial partial criteria are reported. This is supplied dated curation derived from the OWR record, not an independently verified assessment; time-sensitive status and cited claims require independent verification.
\item The record carries the supplied statement-status label \texttt{corrected\_verified}, with an explanation that mathematical content was preserved while notation, typography, omitted definitions, or duplicate clauses were repaired. This label does not establish independent mathematical or source verification of this draft.
\end{itemize}
\end{document}
